\section{Optical Simulation}
\label{sec:optical}

\subsection{The Placement Problem}

Lens artifacts belong to the taking lens and therefore occur \emph{before} the
film is exposed. Diffusion filters, however, may sit on the taking lens or may
be a property of the projection and viewing chain, and bloom in a projected
print is a property of the projector and the screen. \EM{} therefore splits the
optical stage into two passes at different points in the pipeline:

\begin{itemize}
\item \textbf{Pre-exposure optics} (before \eqref{eq:charcurve}): geometric
distortion, vignetting, chromatic aberration, and taking-lens diffusion. These
modify the exposure that forms the latent image, so they interact correctly with
the characteristic curve --- a vignetted corner falls further down the toe and
loses shadow detail, which is what actually happens.
\item \textbf{Post-print optics} (after \eqref{eq:printtodisplay}): bloom and
veiling glare associated with viewing. These act on the rendered image.
\end{itemize}

Applications that apply all optical effects at the end get vignetting
qualitatively wrong: a post-hoc darkening multiplies the output, whereas a real
vignette reduces exposure and therefore changes contrast and color in the
corners, not merely brightness.

\subsection{Geometric Distortion}

The Brown--Conrady model is used, with radial and tangential terms. For a
normalized image coordinate $\mathbf{x} = (x, y)$ measured from the optical
center with $r^2 = x^2 + y^2$:

\begin{align}
\mathbf{x}_d &= \mathbf{x}\left(1 + k_1 r^2 + k_2 r^4 + k_3 r^6\right) \nonumber \\
&\quad + \begin{pmatrix} 2 p_1 x y + p_2 (r^2 + 2x^2) \\ p_1 (r^2 + 2y^2) + 2 p_2 x y \end{pmatrix}.
\label{eq:distortion}
\end{align}

$k_1 < 0$ gives barrel distortion, $k_1 > 0$ pincushion. The tangential terms
$p_1, p_2$ model decentering and are exposed only in the "vintage lens" presets,
where a small asymmetry is a large part of the character.

Distortion is implemented as an inverse warp: for each output pixel, compute the
source coordinate and sample with a Catmull--Rom kernel. Forward warping would
require scatter and produce holes. Note that Apple's lens correction is
\emph{enabled} at decode (Listing~\ref{lst:rawconfig}), so the input is
rectilinear and \eqref{eq:distortion} adds the simulated lens's distortion to a
clean base rather than fighting the phone lens's own.

\subsection{Vignetting}

Two mechanisms are modelled and summed in the illumination domain.

\paragraph{Natural (cos-fourth) falloff} An unavoidable consequence of the
geometry of an ideal lens:

\begin{equation}
V_{\mathrm{nat}}(\theta) = \cos^4 \theta, \qquad \tan\theta = \frac{r\, d}{f},
\end{equation}

where $d$ is half the frame diagonal and $f$ the focal length. Selecting a
shorter simulated focal length therefore automatically produces stronger
falloff, which is correct and which links the vignette control to the format
and lens selection rather than leaving it free-floating.

\paragraph{Mechanical vignetting} Aperture obstruction by the lens barrel,
which is strongly aperture-dependent and vanishes on stopping down:

\begin{equation}
V_{\mathrm{mech}}(r) = 1 - v\,\max\!\left(0, \frac{r - r_0}{1 - r_0}\right)^{\!\eta} \cdot \varsigma_N,
\end{equation}

with $\varsigma_N = \mathrm{clamp}\!\left((N_{\max}/N)^2, 0, 1\right)$ the
aperture dependence, $N$ the f-number, $r_0$ the onset radius, and $\eta$ the
falloff exponent (default $2.2$). The combined operator is

\begin{equation}
\boxed{\;E' = E \cdot V_{\mathrm{nat}}\cdot V_{\mathrm{mech}}\;}
\label{eq:vignette}
\end{equation}

applied to linear exposure, pre-characteristic-curve.

\subsection{Chromatic Aberration}

Lateral chromatic aberration is a wavelength-dependent magnification. It is
modelled as a per-channel radial scale about the optical center:

\begin{equation}
\mathbf{x}_c = \mathbf{x}\left(1 + \varkappa\, \delta_c\, r^2\right),
\qquad \boldsymbol{\delta} = (+1, 0, -1),
\label{eq:ca}
\end{equation}

so the red and blue records are scaled oppositely and the green record is the
reference. $\varkappa$ is the aberration strength. The quadratic $r^2$
dependence rather than a constant scale is important: real lateral CA is
negligible at the center and grows toward the corners, and a constant scale
produces visible fringing on center subjects where there should be none.

Longitudinal chromatic aberration --- the shift of focal plane with wavelength,
which produces the green/magenta fringing on out-of-focus edges --- requires
depth information and is implemented only when a depth map is available from
capture, as a per-channel defocus radius offset.

\subsection{Diffusion Filters}

A diffusion filter (Pro-Mist, Hollywood Black Magic, and similar) contains
particles that scatter a small fraction of the light over a wide angle. The
visible effects are a soft halo around highlights and, characteristically, a
lifting of shadows \emph{adjacent to} highlights --- a local, not global, black
lift. Modelling it as a global blur-and-blend loses that locality and simply
looks out of focus.

The model is a two-term veiling glare:

\begin{equation}
\boxed{\;E' = (1 - w_d) E \;+\; w_d \left[(1-\varrho_d)\, G_{\sigma_1} * E \;+\; \varrho_d\, G_{\sigma_2} * E\right]\;}
\label{eq:diffusion}
\end{equation}

with $\sigma_1$ small (a few pixels, the tight halo), $\sigma_2$ very large (the
broad veil), $\varrho_d$ the split, and $w_d$ the filter strength ($1/8$, $1/4$,
$1/2$, $1$, $2$ presets mapping to $w_d \in \{0.03, 0.06, 0.11, 0.19, 0.30\}$).
Because the operation is in linear exposure, a bright highlight contributes far
more to the veil than a mid-tone, which localizes the effect around highlights
automatically --- the essential behavior, obtained without a threshold.

Diffusion is applied \emph{pre-exposure} when the user selects "taking lens
filter" and \emph{post-print} when they select "projection glare." The two look
meaningfully different: pre-exposure diffusion is compressed by the film's
shoulder and looks restrained; post-print diffusion is not and looks stronger.
Both are offered, and the difference is a small piece of photographic education
delivered by a toggle.

\subsection{Bloom}

Bloom is the post-print counterpart of halation: the spread of very bright areas
in the viewing chain. It uses the same pyramid machinery as
Section~\ref{sec:recursive-blur} with a Gaussian rather than exponential target
and a higher threshold, and is applied to display-linear values. It is disabled
by default, since halation already supplies the film-side glow and stacking both
is a common way to make an image look cheap.

\subsection{Depth-Dependent Effects}

When the capture provides a depth map (available on devices with a LiDAR
scanner or through the dual-camera disparity estimate), three additional effects
become available:

\paragraph{Synthetic defocus} A depth-varying circle of confusion,
\begin{equation}
c(z) = \frac{f^2}{N\,(z_f - f)}\cdot\frac{|z - z_f|}{z},
\end{equation}
with $z_f$ the focus distance. Gathered with a scatter-as-gather approach at
half resolution.

\paragraph{Cat-eye (optical vignetting) bokeh} Off-axis bokeh discs are
truncated by the lens barrel, producing the characteristic lemon shape. The
aperture shape at image radius $r$ is the intersection of two circles of radius
$R$ whose centers are separated by $\varsigma_{\mathrm{cat}}\, r$; the sampling
kernel is masked by that intersection, oriented along the radial direction.

\paragraph{Aperture blade shape} An $n$-gon mask with optional blade curvature,
producing polygonal bokeh for $n \le 9$.

These are marked SHOULD (FR-10) rather than MUST, since depth is not available
for imported files and the quality of phone depth maps at object boundaries is
the limiting factor rather than the optical model.

\subsection{Deferred: Lens Breathing}

Lens breathing --- the change of field of view with focus --- is a temporal
phenomenon with no meaning for a single still frame. It is listed in the
original proposal and is correctly deferred to the video release, where it will
be implemented as a focus-driven scale animation applied before distortion. It
is recorded here so that the omission from v1.0 is deliberate and documented
rather than an oversight.

\subsection{Cost Summary}

\begin{table}[htbp]
\caption{Optical Stage Costs at 12\,MP}
\label{tab:opticalcost}
\begin{center}
\renewcommand{\arraystretch}{1.15}
\begin{tabular}{@{}lll@{}}
\toprule
\textbf{Effect} & \textbf{Passes} & \textbf{Cost} \\
\midrule
Distortion + CA (fused)  & 1 warp    & $1.0\times$ base \\
Vignetting               & fused     & negligible \\
Diffusion                & pyramid   & $1.4\times$ base \\
Bloom                    & pyramid   & $1.4\times$ base \\
Synthetic defocus        & 3 (half)  & $1.1\times$ base \\
\bottomrule
\end{tabular}
\end{center}
\end{table}

Distortion, chromatic aberration, and vignetting are fused into a single warp
kernel: the kernel computes the distorted coordinate once, applies the
per-channel CA scale to it, performs three filtered samples, and multiplies by
the vignette factor. Three separate passes would cost three full-resolution
round trips to memory for no benefit.
